\subsubsection{Ribonucleases in Syn1 and Syn3A}

The ribonuclease complement of \syn~was read from the genome annotation and assigned activities from the characterised \textit{Bacillus subtilis} orthologues.
Both \syn~and \syna~encode the same compact set: the endoribonucleases RNase~III, RNase~Y and RNase~M5, the 5$'\rightarrow$3$'$ exoribonucleases RNase~J1 and RNase~J2, the 3$'\rightarrow$5$'$ exoribonucleases RNase~R and YhaM, and the \textit{trans}-translation factors tmRNA and SmpB (Table~\ref{tbl:rnase_syn1_syn3a})~\cite{durand_rnases_2018, redko_minimal_2013}.
Per-cell abundances were the proteomic copy numbers (Section on proteomics), and substrate signatures were taken from the corresponding \textit{B. subtilis} enzymes.

\subsubsection{RNA-processing endpoint analysis}

Post-transcriptional processing was quantified from the clustered PacBio isoforms.
Isoforms supported by at least ten reads were retained, giving 20{,}885 isoforms.
A per-strand mask of protein-coding gene bodies was built from the \syn~GFF3 annotation, and the 5$'$ and the 3$'$ genomic endpoint of each isoform was labelled intragenic when it fell inside a gene body and intergenic otherwise.
Each isoform was assigned to one of four end-context categories --- unprocessed (both ends intergenic), 5$'$-eroded only, 3$'$-eroded only, or both-eroded --- and the composition was tabulated both by isoform count and by read.
The erosion asymmetry was summarised as the ratio of the read weight carried by 5$'$-intragenic to that carried by 3$'$-intragenic isoforms.
Independently, every annotated start and stop codon fully contained within an isoform was counted, and a reading frame that contained its start codon but not its stop codon was scored as 3$'$-truncated.
For the single-gene panels, the isoforms contained within an operon were coloured by end-context and ordered into the erosion ladder; for the ATP-synthase panel the member isoforms of the two segmented operons were coloured by transcription block.
Analyses were run in Python 3.11 with pandas 2.1, NumPy 1.26 and SciPy 1.11, and figures were drawn with Matplotlib 3.8.

\subsubsection{Homology Comparison of RNase III and RNase Y between B. subtilis and Syn1}

RNase~III and RNase~Y, the endoribonucleases that carry out most intra-operonic cleavage, were compared between \textit{B. subtilis} (NCBI NC\_000964.3) and \syn~to confirm conserved catalytic and RNA-binding architecture.
Domain organisation was assigned with InterPro, and pairwise coverage and identity were taken from the protein alignment.
Both enzymes retained their full domain architecture at 91--98\% coverage and 35--45\% identity (Table~\ref{tbl:rnase_III_Y}).

\begin{table}[H]
% \centering
\caption{RNase III and RNase Y Homology Comparison}
\begin{tabular}{ C{2cm}  C{2cm} C{2cm} C{2cm} C{2cm} C{3cm} }
\toprule
\multicolumn{2}{c}{B. subtilis \textsuperscript{\emph{a}}} & \multicolumn{2}{c}{Syn1} & \multicolumn{2}{c}{Alignment} \\  \hline
Length & Domain \textsuperscript{\emph{b}} & Length & Domain & Coverage \% & Identity \%  \\
\midrule
249 & RNase\_III and dsDNA binding & 232 &  RNase\_III and dsDNA binding & 91 & 45 \\
\hline
520 & RNase\_Y, KH, and HD & 509 & RNase\_Y, KH, and HD & 98 & 35 \\
\bottomrule
\label{tbl:rnase_III_Y}
\end{tabular}

\textsuperscript{\emph{a}} NCBI genome accesion B\_subtilis\_168\_NC\_000964.3 \\
\textsuperscript{\emph{b}} Inferred by InterPro

\end{table}

% \subsubsection{Mapping Endo-RNase Processing Sites from B. subtilis to Syn1}
% TODO (analysis pending, see L2.3/L2.4): map characterised B. subtilis RNase III / RNase Y
% cleavage sites onto Syn1 by homology, and test the 3'-end secondary-structure prediction for the
% 3'->5' exonucleolytic bias (ViennaRNA). Prose to be written once the analysis is done.
